% answers_short.tex —— 2015 年·安徽卷·文科 · 简短答案速查
% 由 tools/make_short_answers.py 自 answers.tex 的 \gksolhead/\gkanswer 提取，
% 并与 2015/2015_数学真题_安徽文_答案要点.md 逐题比对后生成，勿手改（改源 MD 后重跑）。

\begin{gknotice}
\textbf{使用说明}：本节只给最终答案便于快速核对，完整推导见解析版。答案为编写者独立推导，\textbf{非官方标准答案}。
\end{gknotice}

\gkpart{选择题}
\begin{center}\begin{tabular}{ccccccccccc}
\toprule
题号 & 1 & 2 & 3 & 4 & 5 & 6 & 7 & 8 & 9 & 10 \\
\midrule
答案 & C & B & C & A & A & A & B & D & B & A \\
\bottomrule
\end{tabular}\end{center}

\gkpart{填空题}
\begin{itemize}[leftmargin=2.6em]
  \item[\textbf{11.}] \(-1\)
  \item[\textbf{12.}] \(AC=2\)
  \item[\textbf{13.}] \(27\)
  \item[\textbf{14.}] \(-\dfrac{1}{2}\)
  \item[\textbf{15.}] ①④⑤。
\end{itemize}

\gkpart{解答题}
\begin{itemize}[leftmargin=2.6em]
  \item[\textbf{16.}] (1) \(T=\pi\)；(2) 最大值 \(1+\sqrt{2}\)（在 \(x=\dfrac{\pi}{8}\) 处取得），最小值 \(0\)（在 \(x=\dfrac{\pi}{2}\) 处取得）
  \item[\textbf{17.}] (1) \(a=0.006\)；(2) \(0.4\)；(3) \(\dfrac{1}{10}\)
  \item[\textbf{18.}] (1) \(a_n=2^{n-1}\)；(2) \(T_n=1-\dfrac{1}{2^{n+1}-1}\)
  \item[\textbf{19.}] (1) \(\dfrac{\sqrt{3}}{6}\)；(2) 存在（唯一）点 \(M\)，使得 \(AC\perp BM\)；此时 \(\dfrac{PM}{MC}=\dfrac{1}{3}\)
  \item[\textbf{20.}] (1) \(e=\dfrac{2\sqrt{5}}{5}\)；(2) 成立（\(\overrightarrow{MN}\cdot\overrightarrow{AB}=\dfrac{a^2-5b^2}{6}=0\)，由 \(a^2=5b^2\) 得垂直）
  \item[\textbf{21.}] (1) 定义域为 \(\{x\mid x\neq-r\}\)；\(f(x)\) 在 \((-\infty,-r)\) 与 \((r,+\infty)\) 上单调递减，在 \((-r,r)\) 上单调递增；(2) 在 \((0,+\infty)\) 内存在极大值 \(100\)（在 \(x=r\) 处取得），无极小值
\end{itemize}
